\documentclass[11pt,a4paper]{ctexart}
\usepackage{geometry}
\geometry{margin=2.5cm}
\usepackage{listings}
\usepackage{xcolor}
\usepackage{graphicx}
\usepackage{float}
\usepackage{enumitem}

\lstset{
    language=bash,
    basicstyle=\ttfamily\small,
    frame=single,
    breaklines=true,
    showstringspaces=false,
    columns=fullflexible
}

\title{系统开发工具实验报告3}
\author{姓名：李博文 \quad 学号：25020007057}
\date{}

\begin{document}
\maketitle

\section*{第九题}
\begin{lstlisting}
#1. 创建目录结构
mkdir -p q09/src/greetlab
cd q09

#2 新建 __init__.py
touch src/greetlab/__init__.py

#3 写 cli.py
cat > src/greetlab/cli.py <<'EOF'
import argparse

def main():
    p = argparse.ArgumentParser()
    p.add_argument("--name", required=True)
    a = p.parse_args()
    print(f"Hello, {a.name}!")
EOF

#4. 写 pyproject.toml
cat > pyproject.toml <<'EOF'
[build-system]
requires = ["setuptools>=68"]
build-backend = "setuptools.build_meta"

[project]
name = "greetlab-25020007057"
version = "0.1.0"

[project.scripts]
sdt-greet = "greetlab.cli:main"
EOF

#5. 创建打包用虚拟环境
python3 -m venv buildenv
#激活环境
source buildenv/bin/activate

#6 在虚拟环境内安装build
pip install build

#7 执行打包，生成wheel
python -m build

#8 退出打包环境
deactivate

#9 创建干净虚拟环境 testenv
python3 -m venv testenv
source testenv/bin/activate

#10 在 testenv 里，安装打包好的 wheel
pip install dist/*.whl

#11 切换到 q09 文件夹外面测试
cd ../
#执行命令
sdt-greet --name 25020007057

#12 测试完毕，退出虚拟环境
deactivate
\end{lstlisting}
\begin{figure}[H]
    \centering
    \includegraphics[width=0.95\textwidth]{q09.png}
    \caption{第九题运行截图}
\end{figure}

\section*{第十题}
\begin{lstlisting}
#1. 复制整个 q09 目录到 q10
cp -r ~/q09 ~/q10
cd ~/q10

#2. 写入测试文件 test_greetlab.py
cat > test_greetlab.py <<'EOF'
import pytest
from greetlab.cli import main
import sys

def test_blank_name():
    sys.argv = ["cli.py", "--name", "   "]
    with pytest.raises(SystemExit) as excinfo:
        main()
    assert excinfo.value.code == 2
EOF

#3 运行测试，确认失败
#加PYTHONPATH=src 解决找不到模块的问题
PYTHONPATH=src pytest test_greetlab.py -v

#4 替换 cli.py
cat > src/greetlab/cli.py <<'EOF'
import argparse

def main():
    p = argparse.ArgumentParser()
    p.add_argument("--name", required=True)
    a = p.parse_args()
    if len(a.name.strip()) == 0:
        raise SystemExit(2)
    print(f"Hello, {a.name}!")
EOF

#5 人工查看 diff
git diff src/greetlab/cli.py

#6 再次跑测试
PYTHONPATH=src pytest test_greetlab.py -v

#7 创建 ai_log.md（最多 5 行）
cat > ai_log.md <<'EOF'
1.提示：name仅空白时main抛出SystemExit(2)，仅修改cli.py，测试命令PYTHONPATH=src pytest
2.AI改动：增加strip()校验，空白则raise SystemExit(2)
3.人工diff检查：仅新增判断，无多余修改
4.重跑pytest，测试用例通过
EOF
\end{lstlisting}
\begin{figure}[H]
    \centering
    \includegraphics[width=0.95\textwidth]{q10.png}
    \caption{第十题运行截图}
\end{figure}

\section*{第十一题}
\begin{lstlisting}
#1 在 q10 目录新建 communication.md
cat > communication.md <<'EOF'
## Issue

环境：Windows（待确认Python版本）

复现命令：sdt-greet --name " "

实际结果：程序输出 Hello, !，退出码0

期望结果：检测到仅含空白字符的name，以退出码2终止程序

描述：传入仅空格的name参数时，程序未校验空白输入，仍正常打印问候。

## 提交信息

fix: 增加name空白输入校验

问题：当--name仅传入空白字符，程序不做校验，正常输出问候并返回0。

方案：对name执行strip()，去除空白后为空则raise SystemExit(2)。

## 评审意见

Blocking：代码缺少空白参数校验。

风险：用户传入全空白name时程序不会报错，产生无效输出，不符合需求。

建议动作：增加输入校验，当name去除空白为空时，返回退出码2。
EOF

#2 查看文件确认
cat communication.md
\end{lstlisting}
\begin{figure}[H]
    \centering
    \includegraphics[width=0.85\textwidth]{q11.png}
    \caption{第十一题运行截图}
\end{figure}

\section*{第十二题}
\begin{lstlisting}
#1. 创建 q12 目录并进入
mkdir ~/q12
cd ~/q12

#2. 写入完整 train.py（补全两处 TODO）
cat > train.py <<'EOF'
import torch
from torch import nn

torch.manual_seed(20260907)
x = torch.linspace(-1, 1, 101).reshape(-1, 1)
y = 3 * x - 1

model = nn.Linear(1, 1)
loss_fn = nn.MSELoss()
opt = torch.optim.SGD(model.parameters(), lr=0.1)

for _ in range(200):
    pred = model(x)
    loss = loss_fn(pred, y)
    # TODO：清空梯度、反向传播、更新参数
    opt.zero_grad()
    loss.backward()
    opt.step()

# TODO：进入评估模式并在no_grad中打印最终损失、weight和bias
model.eval()
with torch.no_grad():
    pred_final = model(x)
    loss_final = loss_fn(pred_final, y)
    w = model.weight.item()
    b = model.bias.item()
    print(f"Final loss: {loss_final:.6f}")
    print(f"Weight: {w:.4f}")
    print(f"Bias: {b:.4f}")
EOF

#3 创建虚拟环境 venv（在 q12 目录）
python3 -m venv venv

#4 激活虚拟环境
source venv/bin/activate

#5 在虚拟环境里安装 CPU 版 PyTorch
pip install torch --index-url https://download.pytorch.org/whl/cpu

#6 运行训练脚本
python train.py

#7 退出虚拟环境
deactivate
\end{lstlisting}
\begin{figure}[H]
    \centering
    \includegraphics[width=0.95\textwidth]{q12.png}
    \caption{第十二题运行截图}
\end{figure}

\section*{练习}
\subsection*{练习1}
将环境保存到一个文件，创建一个venv，激活它，再激活到另一个文件，然后。环境发生了什么变化？为什么贝壳更喜欢文夫？（提示：观察激活前后。）跑一跑，推理一下停用bash功能到底在做什么。printenv\quad printenv\quad diff before.txt after.txt\quad \$PATH\quad which deactivate
\begin{figure}[H]
    \centering
    \includegraphics[width=0.95\textwidth]{c1.png}
    \caption{练习1截图}
\end{figure}
\begin{enumerate}[label=\arabic*.]
\item 修改 PATH：把 venv/bin 放到 PATH 最前面。 之后执行python、pip，优先调用 venv 里的版本，而不是系统 python。
\item 新增环境变量 VIRTUAL\_ENV，值就是虚拟环境文件夹路径。
\item 定义 shell 函数 deactivate（加载进当前 bash shell 内存）。
\item 修改命令提示符PS1，前面加上(venv)标记，提醒你当前在虚拟环境。
\end{enumerate}

子 shell 里面修改 PATH、定义函数，只影响子 shell；子 shell 退出后，所有修改全部丢弃，父 shell 完全不受影响，环境不会激活。

source activate（等价 . activate）：在当前 shell 进程内直接执行脚本，直接修改当前 shell 内存里的变量、注册 deactivate 函数，环境生效。

\begin{lstlisting}
printenv > before.txt
python3 -m venv venv
source venv/bin/activate
printenv > after.txt
diff before.txt after.txt
echo $PATH
which deactivate
deactivate
\end{lstlisting}

\subsection*{练习2}
创建一个带有 pyproject.toml 的 Python 包，并在虚拟环境中安装。创建一个锁文件并检查它。
\begin{figure}[H]
    \centering
    \includegraphics[width=0.85\textwidth]{c2.png}
    \caption{练习2截图}
\end{figure}

\subsection*{练习3}
安装 Docker，用它用 docker compose 本地搭建缺课课程网站。
\begin{figure}[H]
    \centering
    \includegraphics[width=0.85\textwidth]{c3.png}
    \caption{练习3截图}
\end{figure}

\subsection*{练习4}
为一个简单的Python应用写一个Docker文件。然后写一个docker-compose.yml，运行你的应用和一个Redis缓存。
\begin{figure}[H]
    \centering
    \includegraphics[width=0.85\textwidth]{c4.png}
    \caption{练习4截图}
\end{figure}

\subsection*{练习5}
发布一个Python包到TestPyPI（除非值得分享，否则不要发布到真正的PyPI包！）。然后用该包构建一个 Docker 镜像，推送到 ghcr.io。
\begin{figure}[H]
    \centering
    \includegraphics[width=0.85\textwidth]{c5.png}
    \caption{练习5截图}
\end{figure}

\subsection*{练习6}
为你选择的编程智能体创建并测试一个 AGENTS.md（或等效文件，如 CLAUDE.md）、一个 Skill（如 Claude Code 的 Skill 或 Codex 的 Skill），以及一个子智能体（如 Claude Code 的子智能体）。思考什么情况下会选择用其中某一种而不是另一种。注意，你选的智能体可能不支持其中某些功能；可以跳过或换一个支持的智能体试试。

答：

\subsubsection*{1. CLAUDE.md}
\begin{lstlisting}
cat > CLAUDE.md <<'EOF'
本项目是Python课程设计项目，使用ruff、pytest，遵循PEP8规范。

1. 所有代码修改必须先写测试，TDD风格。
2. 修改代码前，运行代码质量skill校验。
3. 输出优先简短，答辩场景提供精简口述稿。
4. 禁止自动删除文件；删除操作必须人工确认。
5. 遇到复杂重构任务，委托给 refactor 子智能体，不要自己直接大面积改代码。

- Skill: code_quality（ruff + pytest）
- Sub-agent: refactor_agent（代码重构专用）
EOF
\end{lstlisting}

\subsubsection*{2. Skill}
\begin{lstlisting}
mkdir -p skills
cat > skills/code_quality.skill <<'EOF'
# Skill: code_quality
## 描述
运行 ruff check、ruff format --check、pytest，输出质量报告，用于Python代码校验。

## 触发指令
当用户要求：检查代码质量 / 运行静态检查 / 跑测试

## 执行步骤
1. ruff check src/
2. ruff format --check src/
3. pytest -v

## 返回格式
汇总：通过/失败，列出错误，给出修复建议。
EOF
\end{lstlisting}

\subsubsection*{3 子智能体}
\begin{lstlisting}
mkdir -p subagents
cat > subagents/refactor_agent.md <<'EOF'
# Sub-agent: refactor_agent
角色：Python代码重构专家
任务：接收代码片段，做安全重构，保持功能不变。

约束：
1. 禁止修改业务逻辑，只优化结构、命名、消除重复代码。
2. 重构完成后，输出diff，并且必须给出测试命令验证。
3. 输出精简，给出答辩口述要点。

返回格式：
1. 改动diff
2. 验证命令
3. 简短解释
EOF
\end{lstlisting}

什么时候选择哪一种
\begin{enumerate}[label=\arabic*.]
\item CLAUDE.md（AGENTS.md） 选用：项目全局长期规则。整个仓库所有对话都生效，定义项目背景、通用约束、可用资源。
\item Skill 选用：轻量可复用工具 / 固定流程，只是增加工具能力，没有独立身份。
\item 子智能体 Sub-agent 选用：独立专项任务，需要独立推理。子智能体拥有自己角色、目标、思考约束。
\end{enumerate}

\subsection*{练习7}
浏览一个知名项目（例如 Redis 或 curl）的源码。 找出讲义中提到的几种注释类型：有用的 TODO、指向外部文档的引用、解释为何避开某种做法的“why not”注释、或血泪教训。 如果这些注释不存在，会失去什么？
\begin{enumerate}[label=\arabic*.]
\item TODO 注释缺失 丢失项目遗留的已知缺陷 / 优化机会。后续重构时，不知道代码的局限，可能误把临时方案当成永久设计。
\item 外部文档引用缺失 遇到陌生系统 API / 协议时，缺少直接资料入口。维护者需要额外大量时间检索对应文档，提高理解成本。
\item why-not（为什么不这么做）注释缺失【考试高频】 只看到代码 “做了什么”，看不到设计权衡。维护者很容易凭直觉，选用被评估否决的方案，引入平台兼容、性能或者稳定性 bug。
\item 血泪教训注释缺失 丢失历史 bug、安全事故的上下文。后续重构时，很容易误删保护性代码，旧漏洞复现。
\end{enumerate}

\subsection*{练习8}
找一个你熟悉项目里已经合并的 PR，其中包含实质性的评审意见（不只是“LGTM”）。 通读整个评审。 所有评论同样高效吗？ 如果你是 PR 作者，收到这些评论的体验会如何？

答：

不是，效率分等级

高效评论：直接指出并发逻辑漏洞，说明风险在哪，给出修改思路。作者马上理解问题，知道怎么修复。要求补充回归测试，明确目的：防止以后代码改动把 bug 带回来。目标清晰。

中等效率：规范类，有价值，但优先级低于 bug 修复 是一致性规范问题，不影响程序正确性，但维护长期项目很重要。

偏讨论型：不是 bug，是权衡讨论 不是必须修改。评审人提出性能疑问，需要作者分析、解释。这条是讨论，不是强制修改。

作为 PR 作者，收到这种具体化、有理有据的评审，体验是正向的，能帮助写出更稳健代码；反之模糊的指责会打击积极性。

\subsection*{练习9}
用编程智能体完成与代码质量一讲中 Markdown 无序列表正则表达式练习相同的目标。它是通过直接编辑文件来完成的吗？智能体直接编辑文件有什么缺点和局限？想办法提示智能体不通过直接编辑文件来完成。提示：让智能体使用第一讲中提到的某个命令行工具。

答：默认情况下，编程智能体会直接读取并修改本地文件。 直接编辑文件有不少问题：一旦 AI 逻辑出错，会覆盖损坏原始文件；无法提前预览修改结果；而且这个修改过程很难复现。 所以我写提示词约束智能体，不让它直接操作文件，改用第一讲学的 sed 流编辑器。智能体只输出 sed 命令，我们先预览效果，确认没问题再执行写入，更安全、可复现。

\subsection*{练习10}
打开 Stack Overflow，找一个你熟悉技术领域内获得高票回答的问题。 再找一个被关闭或大量被点踩的问题。 按照讲义的建议对比它们：能否预见哪个问题更可能得到好答案？

答：我们完全可以提前预判哪个更容易收获高质量回答。 优质问题：范围聚焦、提供最小可复现代码、说明已尝试的排查步骤，明确预期结果。回答者可以快速理解问题，给出精准答案。 低分 / 被关闭问题：描述模糊，缺少报错与最小样例，直接要求别人完成全部工作。回答者缺少信息，很难有效帮助，容易被投票否决或者直接关闭。

\section*{感想}
做完 9 到 12 题的实验，我亲手体验了从源码构建 Python 包、测试驱动修复、规范协作文档编写，还有 PyTorch 训练循环的整套流程。第 9 题需要搭建项目目录，编写 pyproject.toml，构建 wheel 安装包。动手操作后，我才理解 Python 项目打包的规范，知道 wheel 包是用来方便分发安装的。第 10 题是 TDD 测试驱动修复，先写会失败的测试用例，再修改代码。这一步让我明白，不管是自己改代码还是借助工具修改，都必须人工检查改动，不能直接照搬，要删掉无关的修改，保证改动只针对问题本身。第 11 题练习写规范的 issue、提交信息和代码评审意见。之前写描述总是很随意，做完题目才感受到，清晰准确的文字能减少团队沟通的误会。第 12 题补全 PyTorch 线性回归的训练循环，弄懂梯度清零、反向传播、更新参数的顺序，分清训练模式和评估模式。这一系列练习让我意识到，工程开发不单单是写出能运行的代码。打包、测试、规范沟通，都是保障软件稳定、方便后续维护必不可少的部分。

\end{document}